\documentclass[aspectratio=169]{beamer}
\usepackage[T1]{fontenc}
\usepackage[utf8]{inputenc}
\usepackage{lmodern}
\usepackage{mastering_ir_derivatives_beamer_1.0}
\usepackage{booktabs}
\usepackage{array}
\usepackage{listings}

% Retain the supplied visual system without changing the authored slide count.
\beamerdefaultoverlayspecification{<*>}
\AtBeginSection[]{}
\AtBeginSubsection[]{}

\hypersetup{
    colorlinks=true,
    linkcolor=blue,
    urlcolor=blue
}

\lstset{
    basicstyle=\ttfamily\scriptsize,
    breaklines=true,
    columns=fullflexible
}

\newcommand{\product}{TGIR QuantLib Tools}
\newcommand{\code}[1]{\texttt{#1}}
\newcommand{\docdate}{2026-04-06}


\usepackage{amsmath,amssymb,mathtools}

\definecolor{tgirnavy}{HTML}{17324D}
\definecolor{tgirblue}{HTML}{2D6A9F}
\definecolor{tgirlight}{HTML}{EAF1F7}
\definecolor{tgirgreen}{HTML}{1F7A5A}
\setbeamercolor{structure}{fg=tgirblue}
\setbeamercolor{frametitle}{fg=tgirnavy}
\setbeamercolor{block title}{bg=tgirnavy,fg=white}
\setbeamercolor{block body}{bg=tgirlight,fg=black}
\setbeamercolor{alerted text}{fg=tgirgreen}

\newcommand{\Qd}{\mathbb{Q}^{d}}
\newcommand{\E}{\mathbb{E}}
\newcommand{\F}{\mathcal{F}}
\newcommand{\dd}{\,\mathrm{d}}
\newcommand{\DOM}{\ensuremath{\mathrm{DOM}}}
\newcommand{\FOR}{\ensuremath{\mathrm{FOR}}}

\title[Callable Bermudan XCCY]{Callable Bermudan Cross-Currency Swap}
\subtitle{Quantitative model, simulation, exercise, and validation}
\author{Tim Glauner\\TG Investments and Research, LLC}
\date{August 30, 2026}

\begin{document}

\begin{frame}[plain]
  \titlepage
\end{frame}

\begin{frame}{Executive Summary}
  \begin{columns}[T,onlytextwidth]
    \column{0.55\textwidth}
      \begin{block}{Base model}
        Two Hull-White rate factors, lognormal FX, and Bermudan exercise by
        least-squares Monte Carlo (LSM).
      \end{block}
      \begin{itemize}
        \item Domestic money-market account is the numeraire.
        \item Values and exercise comparisons are in domestic currency.
        \item Deterministic forecast/discount basis is frozen in version 1.
      \end{itemize}
    \column{0.40\textwidth}
      \begin{block}{Core control}
        An independently calculated non-callable XCCY cash-flow PV is the
        primary benchmark and control variate.
      \end{block}
      \begin{alertblock}{Holder perspective}
        Positive cash flows are receipts by the exercise holder.
      \end{alertblock}
  \end{columns}
\end{frame}

\section{Product and Exercise}

\begin{frame}{Product Cash-Flow Anatomy}
  \begin{columns}[T,onlytextwidth]
    \column{0.47\textwidth}
      \begin{itemize}
        \item Domestic and foreign coupon legs.
        \item Optional initial and final notional exchanges.
        \item Exercise dates $t_1<\cdots<t_M$ with explicit settlement timing.
        \item FX quote $S_t$: units of \DOM{} per one unit of \FOR{}.
      \end{itemize}
      \begin{block}{FX conversion at payment}
        \[
          CF_f(T)\longmapsto S_TCF_f(T).
        \]
      \end{block}
    \column{0.48\textwidth}
      \begin{tikzpicture}[
        node distance=5mm,
        box/.style={draw=tgirnavy,rounded corners,fill=tgirlight,
          align=center,text width=2.7cm,minimum height=8mm},
        flow/.style={-{Latex[length=2mm]},thick,draw=tgirblue}
      ]
        \node[box] (common) {Common cash flows\\$A_i$};
        \node[box,below=of common] (exercise) {Exercise settlement\\$J_i$};
        \node[box,below=of exercise] (tail) {Cancellable continuation\\$Y_i$};
        \node[box,right=5mm of exercise] (value) {Decision-date value\\$V_i$};
        \draw[flow] (common.east) -- (value.west);
        \draw[flow] (exercise.east) -- (value.west);
        \draw[flow] (tail.east) -- (value.west);
      \end{tikzpicture}
  \end{columns}
\end{frame}

\begin{frame}{Exercise Decision: Compare Like With Like}
  \begin{block}{Continuation estimate}
    \[
      C_i=\E^{\Qd}\!\left[Y_i\mid\F_{t_i}\right],
      \qquad
      V_i=A_i+\max(J_i,C_i).
    \]
  \end{block}
  \begin{columns}[T,onlytextwidth]
    \column{0.48\textwidth}
      \begin{itemize}
        \item Exercise when $J_i>C_i+\varepsilon_{\mathrm{ex}}$.
        \item Exclude common $A_i$ flows from both alternatives.
        \item Add $A_i$ outside the comparison.
      \end{itemize}
    \column{0.47\textwidth}
      \begin{itemize}
        \item Zero-fee cancellation: $J_i=0$.
        \item Cancel when continuation is negative.
        \item Entry option: $J_i$ is the value of entering the swap plus settlement.
      \end{itemize}
  \end{columns}
  \begin{alertblock}{Implementation consequence}
    The first exercise removes only the cancellable tail; paid, common, and
    delayed cash flows remain.
  \end{alertblock}
\end{frame}

\section{Joint Market Model}

\begin{frame}{Three-Factor Dynamics Under the Domestic Measure}
  \small
  \begin{align*}
    \dd r_d(t) &= [\theta_d(t)-a_dr_d(t)]\dd t+\sigma_d(t)\dd W_d(t),\\
    \dd r_f(t) &= [\theta_f(t)-a_fr_f(t)-\rho_{fS}\sigma_f(t)\sigma_S(t)]\dd t
                  +\sigma_f(t)\dd W_f(t),\\
    \frac{\dd S_t}{S_t} &= [r_d(t)-r_f(t)]\dd t+\sigma_S(t)\dd W_S(t).
  \end{align*}
  \begin{columns}[T,onlytextwidth]
    \column{0.49\textwidth}
      \begin{block}{Quanto adjustment}
        The negative foreign-rate drift adjustment follows from the change to
        the domestic pricing measure for the stated FX quotation.
      \end{block}
    \column{0.46\textwidth}
      \begin{block}{No-arbitrage check}
        \[
          S_0P_f^{*}(0,T)=\E^{\Qd}[D_d(0,T)S_T].
        \]
      \end{block}
  \end{columns}
\end{frame}

\begin{frame}{Curves, Basis, and Correlation}
  \begin{columns}[T,onlytextwidth]
    \column{0.51\textwidth}
      \begin{block}{Deterministic multi-curve mapping}
        For currency $c$, freeze
        \begin{align*}
          B_c(0,T)&=\frac{P_c^F(0,T)}{P_c^D(0,T)},\\
          P_c^F(t,T)&=P_c^D(t,T)\frac{B_c(0,T)}{B_c(0,t)}.
        \end{align*}
      \end{block}
      \begin{itemize}
        \item One stochastic rate factor per currency.
        \item Stochastic tenor or cross-currency basis is out of scope.
      \end{itemize}
    \column{0.44\textwidth}
      \begin{block}{Instantaneous correlation}
        \[
        R=\begin{pmatrix}
          1&\rho_{df}&\rho_{dS}\\
          \rho_{df}&1&\rho_{fS}\\
          \rho_{dS}&\rho_{fS}&1
        \end{pmatrix}
        \]
        \[
          R=R^{\mathsf T},\qquad R\succeq0.
        \]
      \end{block}
      Govern estimates, validate positive semidefiniteness, and stress all
      three correlations independently.
  \end{columns}
\end{frame}

\section{Calibration and Simulation}

\begin{frame}{Calibration Stack}
  \begin{columns}[T,onlytextwidth]
    \column{0.49\textwidth}
      \begin{enumerate}
        \item \textbf{Curves:} bootstrap domestic collateral and foreign
          effective curves; reprice helpers and enforce FX-forward parity.
        \item \textbf{Domestic HW:} calibrate to normal-vol swaptions on the
          exercise diagonal; constrain $a_d$ and regularize $\sigma_d(t)$.
        \item \textbf{Foreign HW:} apply analogous identification controls.
      \end{enumerate}
    \column{0.46\textwidth}
      \begin{enumerate}
        \setcounter{enumi}{3}
        \item \textbf{FX diffusion:} fit $\sigma_S(t)$ to selected vanillas in
          the full hybrid model.
        \item \textbf{Correlation:} combine governed historical/implied
          estimates with PSD checks and stresses.
      \end{enumerate}
      \begin{block}{Persist the evidence}
        Parameters, helper prices, relative errors, optimizer status, and
        curve/vol/correlation snapshot IDs.
      \end{block}
  \end{columns}
\end{frame}

\begin{frame}{Monte Carlo Event Grid and Path Evolution}
  \begin{columns}[T,onlytextwidth]
    \column{0.49\textwidth}
      \begin{block}{Event grid contains every}
        Fixing, accrual, payment, notional exchange, decision, effective, and
        settlement date, plus capped intermediate steps.
      \end{block}
      \begin{itemize}
        \item Generate correlated normals from a factorization of $R$.
        \item Use exact conditional Gaussian Hull-White transitions when available.
        \item Integrate domestic short rates consistently for discounting.
      \end{itemize}
    \column{0.46\textwidth}
      \begin{block}{Log-FX evolution}
        \[
          \begin{aligned}
          \dd\log S_t
            &=\left[r_d(t)-r_f(t)-\tfrac12\sigma_S^2(t)\right]\dd t\\
            &\quad+\sigma_S(t)\dd W_S(t).
          \end{aligned}
        \]
      \end{block}
      \begin{itemize}
        \item Apply analytic conditional bond values where possible.
        \item Start with antithetics; validate low-discrepancy variants separately.
      \end{itemize}
  \end{columns}
\end{frame}

\begin{frame}{Two-Pass Least-Squares Monte Carlo}
  \begin{tikzpicture}[
    node distance=7mm and 8mm,
    box/.style={draw=tgirnavy,rounded corners,fill=tgirlight,
      align=center,text width=3.6cm,minimum height=11mm},
    flow/.style={-{Latex[length=2mm]},thick,draw=tgirblue}
  ]
    \node[box] (train) {\textbf{Training paths}\\work backward from $t_M$};
    \node[box,right=of train] (regress) {Regress continuation\\on standardized state basis};
    \node[box,right=of regress] (freeze) {Freeze preprocessing,\\basis, and coefficients};
    \node[box,below=of freeze] (price) {\textbf{Independent pricing paths}\\apply policy forward};
    \node[box,left=of price] (cash) {Stop cancellable tail;\\retain common flows};
    \node[box,left=of cash] (npv) {Discount realized flows;\\average NPV and diagnostics};
    \draw[flow] (train) -- (regress);
    \draw[flow] (regress) -- (freeze);
    \draw[flow] (freeze) -- (price);
    \draw[flow] (price) -- (cash);
    \draw[flow] (cash) -- (npv);
  \end{tikzpicture}
  \vspace{0.6em}
  \begin{block}{Regression evidence}
    Record rank, condition number, coefficients, path count, and residual
    diagnostics. Re-train across seeds and compare policies on one common
    pricing sample.
  \end{block}
\end{frame}

\section{Outputs, Risk, and Validation}

\begin{frame}{Decision-Ready Outputs and Risk}
  \begin{columns}[T,onlytextwidth]
    \column{0.49\textwidth}
      \begin{block}{Report}
        Callable and non-callable NPV, embedded option value, Monte Carlo
        confidence interval, exercise probabilities, expected exercise time,
        per-leg PV, and calibration/regression diagnostics.
      \end{block}
      \[
        V_{\mathrm{option}}=V_{\mathrm{callable}}-V_{\mathrm{noncallable}}\ge0
      \]
    \column{0.46\textwidth}
      \begin{block}{Bump-and-revalue with common random numbers}
        \begin{itemize}
          \item Domestic and foreign curve DV01.
          \item Forecast- and cross-currency-basis risk.
          \item Rate vegas and FX delta/gamma/vega.
          \item Three correlations and Hull-White parameter stresses.
        \end{itemize}
      \end{block}
      State whether every bump recalibrates the model and retrains the policy.
  \end{columns}
\end{frame}

\begin{frame}{Minimum Validation Ladder}
  \small
  \begin{enumerate}
    \item No exercise dates reproduce the non-callable XCCY PV.
    \item Zero-volatility and identical-currency limits reproduce deterministic
      and single-currency results.
    \item One exercise date matches an independent European or nested-MC benchmark.
    \item Domestic-discounted domestic assets and FX-converted foreign assets
      are martingales under $\Qd$.
    \item Prices converge across paths, time grids, regression bases, and seeds.
    \item Multi-exercise cases bracket the LSM lower bound with a dual upper
      bound or a trusted nested-MC benchmark.
    \item Differences use paired pathwise or covariance-aware errors under
      common random numbers.
  \end{enumerate}
  \begin{alertblock}{Acceptance principle}
    Model fit alone is insufficient: exercise logic, event ordering, policy
    stability, numerical convergence, and limiting cases must all pass.
  \end{alertblock}
\end{frame}

\begin{frame}{Implementation Boundary and References}
  \begin{columns}[T,onlytextwidth]
    \column{0.52\textwidth}
      \begin{block}{Version 1 delivers}
        A transparent three-factor hybrid reference model with controlled
        calibration, independent LSM policy evaluation, and auditable outputs.
      \end{block}
      \begin{block}{Requires a new model version}
        Stochastic basis, smile dynamics beyond benchmark fitting, additional
        rate factors, collateral switching, or richer exercise rights.
      \end{block}
    \column{0.43\textwidth}
      \textbf{Primary references}
      \begin{itemize}
        \item Longstaff and Schwartz (2001), ``Valuing American Options by Simulation.''
        \item QuantLib documentation: \url{https://www.quantlib.org/docs.shtml}.
      \end{itemize}
      \vspace{0.5em}
      \textbf{Source note}

      This deck summarizes the TGIR callable Bermudan XCCY quantitative model
      article dated August 30, 2026.
  \end{columns}
\end{frame}

\begin{frame}[plain]
  \begin{center}
    \vfill
    {\usebeamerfont{title}\color{tgirnavy}Callable Bermudan XCCY Model\par}
    \vspace{0.75em}
    {\large Questions and discussion\par}
    \vspace{1.5em}
    {\small TG Investments and Research, LLC\par}
    \vfill
  \end{center}
\end{frame}

\end{document}
